\section{Positive norms of the original eventual test-configuration base ideals}

\begin{lemma}\label{lem:model-ideal-positive-on-unit-restriction}
\leanok
\lean{EnvelopeProperty.Berkovich.RelativeSchemeValuedPoint.modelIdealNorm_pos_of_comap_eq_top}
Let $\mathcal X$ be proper over an affine valued base whose coefficients
have value at most one. Let $I$ be an original ideal sheaf and $U$ an
original open with $I\mathcal O_U=\mathcal O_U$. At every native valued
point whose underlying scheme point belongs to $U$, the original model
ideal norm of $I$ is strictly positive.
\end{lemma}
\begin{proof}
The support of the pulled-back ideal is the inverse image of the original
support. The unit ideal has empty support, so the scheme point is outside
the original support. The actual model ideal norm vanishes exactly at the
native scheme support. It is therefore nonzero and hence strictly positive.
The valuative model center may still lie in the support; the argument uses
the original native point, not its center.
\end{proof}

\begin{lemma}\label{lem:native-test-configuration-valued-point-unit-open}
\leanok
\lean{EnvelopeProperty.NativeTestConfiguration.TestConfiguration.valuedPoint_mem_unitOpen}
For an actual valuation $v$ on the original polynomial parameter ring
with $v(t)\ne0$, every actual native point extending $v$ lies in the
original noncentral scheme open of a literal test configuration.
\end{lemma}
\begin{proof}
The native point's whole prescribed-base equation identifies the image of
its scheme point with the support of $v$. This support belongs to $D(t)$
exactly when $v(t)\ne0$. Taking the original inverse image gives membership
in the original noncentral open.
\end{proof}

\begin{theorem}\label{thm:native-test-configuration-eventual-positive-base-ideal-norm}
\leanok
\lean{EnvelopeProperty.NativeTestConfiguration.TestConfiguration.exists_eventual_powerBaseIdeal_norm_pos}
Take the original ample fibre line and source line-action/fibre-one data
of the eventual-verticality theorem. Fix any bounded valuation on the
original polynomial parameter ring with nonzero parameter value. There
is one $N$ such that for every $m\ge N$ the model norm of the base ideal
of \emph{all} global sections of the original total line power is
strictly positive at every native point extending that valuation.
\uses{thm:native-test-configuration-eventual-verticality,lem:model-ideal-positive-on-unit-restriction,lem:native-test-configuration-valued-point-unit-open}
\end{theorem}
\begin{proof}
Use the single bound derived from original ampleness in eventual
verticality. For each exponent beyond it, the entire original base ideal
restricts to the unit ideal on the original noncentral open. Every native
point extending the fixed valuation lies in this open by the preceding
lemma. The first lemma gives strict positivity. Source projectivity gives
the properness required by the actual model norm construction.
\end{proof}

\begin{theorem}\label{thm:native-test-configuration-eventual-finite-positive-norm-family}
\leanok
\lean{EnvelopeProperty.NativeTestConfiguration.TestConfiguration.exists_eventual_finite_power_section_norms}
For every exponent beyond one bound, one finite family of actual global
sections of the same original total line power computes its entire base
ideal norm at every original native point. The corresponding maximum of
actual generic-fibre model line norms is strictly positive and at most one.
The family is chosen before the native point; no chosen-family norm
identity, positivity, or verticality certificate is assumed.
\uses{thm:native-test-configuration-eventual-positive-base-ideal-norm}
\end{theorem}
\begin{proof}
The original polynomial base is Noetherian and the original projection
has its actual projective factorization. The whole projective base-ideal
norm theorem therefore gives one finite family whose maximum of actual
generic-fibre section norms equals the all-global-section base ideal norm
at every point. Strict positivity follows from the preceding theorem,
and the bound by one follows from boundedness of the actual ideal norm.
The maximum is a nonnegative real number, and its inclusion in the
extended nonnegative reals preserves these strict and weak inequalities.
\end{proof}
